
\subsection{研究意义}
%%%%%%%%
%\newpage
\vspace{-5pt}

普通引用\cite{test}；上标引用\citess{test}；多篇文章\citess{test,test2,test3}。

有注音的英文：\cite{test}。

参考期刊\cite{test}；
参考图书\cite{test2}；
参考会议\cite{test5}；
参考链接\cite{test4}；
参考文件\cite{test6}。

无序列表的例子：
\begin{itemize}
	\item[-] 第一条；
	\item[-] 第二条。
\end{itemize}

有序列表的例子：
\begin{enumerate}
	\item 第一条；
	\item 第二条。
\end{enumerate}

图片的例子：
\begin{figure}[h]
\centering %图片居中
\includegraphics[width=5cm]{figures/wechatgroup}
\captionsetup{justification=centering} %图题居中
\caption{这是一个很长长长长长长长长长长长长长长长长长长长长长长长长长长长长长长长长长长长长长长长长长长长长长长长长长长的图题。}
\end{figure}


在表格内的第一行设置\verb|\zhkai\ensong\selectfont|，来选择字体，其中\verb|\zhkai\zhsong\enkai\ensong|可以根据需要选择。
\begin{table}[htbp]
	\zhkai\ensong\selectfont%设置表格字体
	\centering  % 显示位置为中间
	\caption{表格}  % 表格标题
	\label{table1}  % 用于索引表格的标签
	%字母的个数对应列数，|代表分割线
	% l代表左对齐，c代表居中，r代表右对齐
	\begin{tabular}{|c|c|c|c|}  
		\hline  % 表格的横线
		& & & \\[-6pt]  %可以避免文字偏上来调整文字与上边界的距离
		第一列&第二列&第三列&第四列 \\  % 表格中的内容，用&分开，\\表示下一行
		\hline
		& & & \\[-6pt]  %可以避免文字偏上 
		0.1&0.2&0.3&0.4 \\
		\hline
	\end{tabular}
\end{table}

对于中文参考文献，bib条目中需要有language = {zh}，参见\cite{test2}。

\begin{REF}
\subsection*{参考文献}
\vspace{-50pt}
\bibliographystyle{gbt7714-nsfc}
\bibliography{ref}%参考文献
\end{REF}

